%%%%%%%%%%%%%%%%%%%%%%%%%%%%%%%%%%%%%%%%%%%%%%%%%%%
\section{Calibration}
% ------------------------------------------------


Calibration is the process of removing the effects of instrumental and atmospheric factors in visibility data before they are combined into an image \citep{Thompson}. The best calibration strategy is to observe a calibration source, an object with well-known properties. Different calibration steps require calibrators with different properties, such as brightness, size, and location close to the science target.





Before calibration begins, contaminated data due to random glitches or bad data values are flagged and removed to make calibration as effective as possible \citep{Fissel2026, Thompson}.



The first step of the calibration process is bandpass calibration, which corrects for frequency-dependent instrumental effects \citep{CASA2012,Ricci2017}. The best calibrator source for this is a very bright quasar with a flat spectrum that may not be located near the science target \citep{Bendo2026}.



The next step is phase and amplitude gain calibration, which corrects time-dependent atmospheric and instrumental variations during the observation \citep{Ricci2017}. A bright, point-like quasar close to the science target is used as a gain calibrator as it is observed frequently throughout the observation \citep{Ricci2017}.







The last step is flux calibration, which scales the relative amplitudes to the correct flux density scale \citep{Ricci2017}. A good flux calibrator is a source with a well-known flux density, usually a solar system object or quasar \citep{Ricci2017}


\SC{comment}

% \quoted{Flux (amplitude) calibration: two steps. 1. Use calibration devices with knwon temperatures (hotload and ambient) to measure System Temperature frequently. 2. Use a source of knwn flux to convert the signal measured at the antenna to a common unit (Janskey, Jy). If the calibrator is resolved, r has sppectral lines, a model for the source is needed. The derived amplitude vs time corrections for the flux calibrator are then applied to the science target.} \citep{Ricci2017}.






%%%%%%%%%%%%%%%%%%%%%%%%%%%%%%%%%%%%%%%%%%%%%%%%%%%

% Ideal calibration sources are to be unresolved since the response to such a source is predictable and should vary only slowly and smoothly with time, and precise details of its visibility are not required; a strong calibrator with a good SNR in a short time, and the position of the calibrator should be close to the target position  \citep{Thompson}.


% There are teo categories. First, slow changing: antenna position, pointing or gain corrections with elevation, shadowing, can be characterized then applied to all observations \citep{Fissel2026}. Secondly fast chaning: variable phaes delay from the atmosphere, flcuations in system temperature, phase drifts in the LO, must be monitored during science observations by oberving calibrators \citep{Fissel2026}. 


% \quoted{Bandpass: Correct frequency-dependent telescope response} \citep{Ricci2017}.


% \quoted{Visibilities are calibrated by determining the complex gains (amplitude and phase), the frequency response (bandpass) and flux scale for each antena} \citep{Ricci2017}.


% Typically, you want the calibration source to be close to your target as the antennas have to point to the calibration source multiple times through observations \add{cite?}. 
%%%%%%%%%%%%%%%%%%%%%%%%%%%%%%%%%%%%%%%%%%%%%%%%%%%


% \quoted{first pass is atmospheric correction from water vapour radiometers readings. final correction from gain calibrator (point source near the target) is observed every few minuites throughout the observation} \citep{Ricci2017}.



% \quoted{Phase calibratin: This provides a moel of atmospheric phase change along the line of sight to the science target that can be comesated for in the data } \citep{Ricci2017}.

% \quoted{Phase and amplitude gain calibratoron: remove effects of time-varying phases/amplitude due to atmospheric or instumental variations } \citep{Ricci2017}.
% \quoted{Atmospheric phase correction: variations in the amount of predictable water vapour cause fluctuations, patches of air with different water vapour (and hence index of refraction) affect the incoming wave front differently} \citep{Ricci2017}.


